\documentclass[11pt,a4paper]{article}

% ---- Packages ----
\usepackage[T1]{fontenc}
\usepackage[utf8]{inputenc}
\usepackage{palatino}
\usepackage{mathpazo}
\usepackage{microtype}
\usepackage{geometry}
\geometry{margin=1in, headheight=14pt}
\usepackage{amsmath,amssymb,amsthm}
\usepackage{booktabs}
\usepackage{hyperref}
\hypersetup{
  hidelinks,
  pdftitle={Certified Quantised Attention with Runtime Precision Escalation},
  pdfauthor={Dean Calver},
  pdfsubject={KV cache quantisation with formal error bounds and FP16 fallback},
  pdfkeywords={KV cache, quantisation, attention, LLM inference, error bounds}
}
\usepackage{fancyhdr}
\usepackage[table]{xcolor}
\usepackage{parskip}
\usepackage{enumitem}
\usepackage{titlesec}
\usepackage{graphicx}
\usepackage{float}
\usepackage{tikz}
\usetikzlibrary{arrows.meta,positioning,fit,backgrounds,calc}

% ---- Page style ----
\pagestyle{fancy}
\fancyhf{}
\rhead{\textit{\small Certified Quantised Attention with Runtime Precision Escalation}}
\cfoot{\thepage}
\renewcommand{\headrulewidth}{0.4pt}

% ---- Theorem environments ----
\newtheorem{theorem}{Theorem}[section]
\newtheorem{corollary}[theorem]{Corollary}
\newtheorem{lemma}[theorem]{Lemma}
\theoremstyle{definition}
\newtheorem{definition}[theorem]{Definition}
\theoremstyle{remark}
\newtheorem{remark}[theorem]{Remark}

% ---- Spacing ----
\setlength{\parskip}{6pt}
\setlength{\parindent}{0pt}
\titlespacing*{\section}{0pt}{18pt}{8pt}
\titlespacing*{\subsection}{0pt}{12pt}{4pt}

% ---- Notation ----
\newcommand{\Kb}{\mathbf{K}}
\newcommand{\Qb}{\mathbf{Q}}
\newcommand{\Vb}{\mathbf{V}}
\newcommand{\R}{\mathbb{R}}
\newcommand{\norm}[1]{\lVert #1 \rVert}
% breakable long identifier
\newcommand{\torchsdpa}{\texttt{torch.scaled\_dot\_\hspace{0pt}product\_attention}}

% =============================================================
\begin{document}

% ---- Title ----
\begin{center}
  {\LARGE\bfseries Certified Quantised Attention}\\[4pt]
  {\LARGE\bfseries with Runtime Precision Escalation}\\[10pt]
  {\large\itshape Provable Error Bounds for Compressed KV Caches with Unconditional FP16 Fallback}\\[8pt]
  Dean Calver\\[4pt]
  {\small\color{gray} April 2026}
\end{center}

\hrule
\vspace{10pt}

% ---- Abstract ----
\begin{abstract}
KV cache quantisation reduces the memory cost of long-context LLM inference, but introduces an approximation error whose magnitude is difficult to characterise in advance.  Existing systems treat this as an acceptable tradeoff: compress aggressively, hope the model is robust.

We present a tiered KV cache that stores INT8 keys and INT4 values in GPU memory while retaining FP16 originals in system RAM\@.  A two-term error decomposition provides independent, per-head, per-step bounds on key and value compression error.  At runtime, the system monitors these bounds and adaptively selects precision per block: blocks with acceptably small error execute on compressed data; blocks that exceed the bound have their FP16 originals paged in.  The worst case is not silent quality degradation---it is the Dense baseline's own \torchsdpa{} code path with unmodified FP16 KV\@.

On LLaMA~3.1-8B across 4K--32K contexts, the certified system matches dense perplexity within noise on PG-19 and is statistically indistinguishable from dense on NIAH (100 paired trials, 8K).  The system runs at ${\sim}2.2\times$ dense latency; effective VRAM savings require 128K+ context due to GQA working-set saturation at shorter lengths.  No model retraining or dataset-specific calibration is required.
\end{abstract}

\vspace{8pt}
\hrule
\vspace{10pt}

\tableofcontents

\newpage

% NOTE: Full paper body supplied by the user in chat should be restored here.
% This placeholder keeps the repository scaffold valid while the extraction work
% below identifies the minimal code and result files that back the paper.

\section{Introduction}
\label{sec:intro}

Autoregressive LLM decoding at long context lengths is dominated by the cost of reading the key-value (KV) cache from GPU memory.  Each decode step fetches cached keys and values for every attention head, a cost linear in sequence length.  At 128K context on a model with 32 attention heads and head dimension 128, the KV cache alone occupies several gigabytes and consumes the majority of available memory bandwidth during decode.

KV cache quantisation addresses this by storing keys and values in reduced-precision formats---typically INT8, INT4, or lower.  The savings are substantial: an INT8 cache halves memory usage, and INT4 quarters it.  But quantisation introduces error, and the relationship between quantisation parameters and output quality is complex.

This paper provides that guarantee through tiered storage with unconditional fallback, formal error bounds with runtime monitoring, and adaptive per-block precision selection.

\section{Results}
\label{sec:results}

The full results tables from the supplied paper should be restored here when the paper is normalised from the chat transcript.

\end{document}
